% 计数测度c与Lebesgue测度lambda在[0,1]上：c给每个单点质量1（示意画出若干等高针），lambda给单点0、给[0,1]质量1；逐点比值h(x)=lambda({x})/c({x})=0，积分为0而lambda([0,1])=1。横向[0,1]对应7cm。
\begin{tikzpicture}[
  lab/.style={font=\figurefont\fontsize{8}{10}\selectfont},
  note/.style={font=\FigureNoteFont\fontsize{7.5}{9}\selectfont,text=FigureMuted}]
\def\L{7}
% counting measure row
\begin{scope}[yshift=2.3cm]
  \draw[FigureStroke,line width=.7pt] (0,0)--(\L,0);
  \foreach \i in {0,...,40}{\draw[FigureBlue,line width=.6pt] ({\i*\L/40},0)--({\i*\L/40},0.7); \fill[FigureBlue] ({\i*\L/40},0.7) circle (0.9pt);}
  \node[lab,anchor=east] at (-0.2,0.35) {计数测度$c$};
  \node[note,anchor=west] at (\L+0.2,0.55) {$c(\{x\})=1$};
  \node[note,anchor=west] at (\L+0.2,0.15) {$c([0,1])=\infty$};
\end{scope}
% Lebesgue row
\begin{scope}
  \fill[FigureGreenFill] (0,0) rectangle (\L,0.7);
  \draw[FigureGreen,line width=1.1pt] (0,0.7)--(\L,0.7);
  \draw[FigureStroke,line width=.7pt] (0,0)--(\L,0);
  \node[lab,anchor=east] at (-0.2,0.35) {长度$\lambda$};
  \node[note,anchor=west] at (\L+0.2,0.55) {$\lambda(\{x\})=0$};
  \node[note,anchor=west] at (\L+0.2,0.15) {$\lambda([0,1])=1$};
  \node[lab,below] at (0,0) {$0$}; \node[lab,below] at (\L,0) {$1$};
\end{scope}
\draw[FigureRed,line width=.7pt,->] (0.6*\L,2.25)--(0.6*\L,0.8);
\node[note,anchor=west,text=FigureRed] at (0.6*\L+0.1,1.55) {逐点比值$h(x)=0/1=0$};
\node[note,anchor=east] at (0.55*\L,1.55) {$\int_{[0,1]}h\,\mathrm dc=0\ne1$};
\end{tikzpicture}
